% 06-hardware.tex

\section{Hardware Adaptation}\label{sec:hardware}

Longitude and Augur are designs; the hardware profiling interfaces they depend on (CUPTI, NVML, rocprofiler) are standard vendor APIs.
\code{CostModel.h} implements the hardware specification and cost model structures.

\subsection{Hardware Profiling}

CUPTI (NVIDIA) and rocprofiler (AMD) provide per-kernel counters: SM/CU utilization, memory bandwidth saturation, achieved occupancy, tensor core utilization, register spill count, cache hit rates, warp execution efficiency, pipeline stall reasons.
NVML and rocm-smi provide device-level metrics: clocks, power, temperature, ECC errors, throttle events, usable memory.
Crucible abstracts these behind a vendor-agnostic profiling interface.

\subsection{Longitude: Startup Calibration}

Longitude runs at startup and on topology change (Relay failure, new Relay, sustained drift detected by Augur).
It populates the \code{HardwareProfile} struct (\code{CostModel.h}) with measured values, replacing the nominal presets (B200, H100, MI300X, A100).

\textbf{Phase 1: GPU profiling.}
Per-GPU, in parallel: GEMM benchmark measures sustained compute throughput (a thermally throttling GPU delivers less than spec).
Streaming copy measures actual HBM bandwidth.
Clock readings under load reveal true boost frequency.
NVML/rocm-smi report power, temperature, ECC errors, usable memory after driver overhead.
Output: a calibrated \code{HardwareProfile} per device with measured values for all fields (\code{num\_sms}, \code{smem\_per\_sm}, \code{hbm\_bw}, \code{peak\_fp16}, etc.).

\textbf{Phase 2: Network probing.}
All-pairs: ping-pong latency, unidirectional and bidirectional bandwidth, topology detection (NVSwitch, NVLink, PCIe, InfiniBand, RoCE, TCP).
Output: weighted network graph of actual measured connectivity.

\textbf{Phase 3: Kernel calibration.}
Representative kernels benchmarked at several shapes.
Ratio of actual measured time to analytical prediction yields a correction factor per kernel class, absorbing cache behavior, instruction scheduling, TLB pressure.
With correction factors, model predictions track measurements within a few percent.

\textbf{Phase 4: Configuration optimization.}
Given measured profiles and topology, formulate system configuration as SMT constraint optimization.
Variables: parallelism degrees (TP, DP, PP, EP, CP), GPU-to-group assignment, communication algorithm per collective per message size, gradient bucket sizes, per-tensor checkpointing decisions, per-op precision assignment, batch size.
Constraints: memory limits, communication feasibility, hardware capabilities.
Objective: minimize predicted iteration time.
Result: complete configuration optimal within the encoded model.

\subsection{Augur: Continuous Monitoring}

Augur runs every iteration, comparing predictions against measurements.

\textbf{Digital twin.}
From the Merkle DAG, kernel cost predictions (roofline model with correction factors), and the memory plan, Augur predicts each iteration.
Per-kernel: execution time, utilization, bottleneck classification (\code{COMPUTE}, \code{MEMORY\_BW}, \code{COMMUNICATION}, \code{BUBBLE}, \code{IMBALANCE}).
Per-iteration: critical path through the DAG, forward/backward/optimizer/communication breakdown.

\textbf{Drift detection.}
Compare predicted vs actual per iteration.
Small sustained deviations update correction factors.
Large sustained deviations trigger diagnosis: thermal throttling (NVML clock drop), hardware degradation (ECC error increase), network congestion (latency spike), workload change (dynamic shapes).
Confirmed hardware change triggers Longitude recalibration.

\textbf{Bottleneck identification.}
Classify dominant bottleneck, rank optimizations by $\text{expected\_speedup} \times \text{confidence}$.
This drives Keeper recommendations.

\subsection{Model Intelligence}

Periodically, Augur computes diagnostics of the Vigil's training dynamics from actual tensor data:

\begin{itemize}
  \item \textbf{Hessian spectrum.} Top eigenvalues via Lanczos iteration (Hessian-vector products, $O(N)$ cost each). Yields smoothness, condition number, convergence rate bounds [Nesterov 1983].
  \item \textbf{Gradient health.} Per-layer gradient norms, signal-to-noise ratio, Jacobian singular values. Detects vanishing/exploding gradients with layer attribution.
  \item \textbf{Representation capacity.} Per-layer effective rank via randomized SVD [Halko et al.\ 2011]. Identifies overparameterization and dead neurons (near-zero variance).
  \item \textbf{Layer redundancy.} CKA [Kornblith et al.\ 2019] between adjacent layers. High CKA $\to$ layers compute nearly identical functions $\to$ pruning candidate.
  \item \textbf{Convergence prediction.} Exponential loss curve fit with confidence intervals. Chinchilla scaling law analysis [Hoffmann et al.\ 2022].
\end{itemize}

Each diagnostic is a standard technique; Crucible's contribution is integrating them into the runtime to drive automated optimization (Section~8).
